% MAIN SNAPSHOT — LSPV/RMV::SYNTHESIS-0020::2026-10-12T02:48:00Z
\documentclass[11pt]{article}
\usepackage[margin=.82in]{geometry}
\usepackage{amsmath,amssymb,amsthm}
\usepackage{booktabs,array,enumitem,xurl,hyperref,microtype}
\hypersetup{colorlinks=true,urlcolor=blue,linkcolor=blue,citecolor=blue}
\setlength{\emergencystretch}{3em}
\newcommand{\EXACT}{\textbf{PROVED/EXACT}}
\newcommand{\KNOWN}{\textbf{KNOWN THEOREM}}
\newcommand{\COMP}{\textbf{COMPUTATIONAL}}
\newcommand{\COND}{\textbf{CONDITIONAL}}
\newcommand{\OPEN}{\textbf{OPEN}}
\newcommand{\RETIRED}{\textbf{FALSIFIED/RETIRED}}
\newcommand{\SAT}{\mathsf{SAT}}
\newcommand{\claim}[1]{\texttt{\detokenize{#1}}}
\newtheorem{theorem}{Theorem}
\newtheorem{corollary}[theorem]{Corollary}
\title{LSPV-5 SYNTHESIS-0020\\
Source specificity is not a computational obstruction}
\author{Universal-core synthesis}
\date{2026-10-12 02:48 UTC}

\begin{document}
\maketitle

\begin{center}
\fbox{\parbox{0.93\linewidth}{\textbf{GLOBAL STATUS: OPEN.}
P versus NP remains open.  The new theorem below is exact but negative: it
separates source distinction from decision hardness.  It supplies no
unrestricted SAT lower bound and no polynomial SAT algorithm.}}
\end{center}

\section{Claim card and synthesis target}

\begin{description}[leftmargin=2.8cm,style=nextline]
\item[Claim ID.] \claim{SOURCE-SPECIFICITY-NO-HARDNESS1}.
\item[Status before run.] \OPEN: FIBERDICH1 showed that an online relation
must fail a same-answer fiber intersection to require source information.
The remaining question was whether that failure itself is progress toward a
universal SAT lower bound.
\item[Exact prefix.] For every set $X$, bit $b:X\to\{0,1\}$, set $Y$, and
map $h:X\to Y$, classify the singleton response relation
\[
 R_h(x,\star)=\{h(x)\}.
\]
Then give an explicit polynomial-time source tag on a same-answer clean-core
pair.
\item[Model.] Deterministic one-way source-removal protocol.  Before the cut
the compiler sees the source; after the cut it must answer one fixed
challenge.  All preprocessing, retained state, response evaluation, and
output are charged.
\item[SAME\_REALIZATION.] The theorem concerns the new source protocol.  Its
cheap tag state is not charged to an unrelated one-bit SAT execution.
\item[Weakest useful result.] Falsify source specificity alone as a
hardness certificate, while preserving the stronger two-way quantitative
compiler/endpoint target.
\end{description}

\section{Exact singleton-relation theorem}

For $a\in b(X)$, write $X_a=\{x:b(x)=a\}$.

\begin{theorem}[\claim{SOURCE-SPECIFICITY-NO-HARDNESS1}]
For $R_h(x,\star)=\{h(x)\}$, the following are equivalent:
\begin{enumerate}[label=(\roman*)]
\item every decision fiber has a common valid response;
\item $h$ is constant on every $X_a$;
\item there exists $f:b(X)\to Y$ with $h=f\circ b$.
\end{enumerate}
If $h$ is not constant on a decision fiber, then the fiber-intersection test
fails, but a source protocol answers exactly by retaining or recomputing
$h(x)$.  In particular, when $h$ is a one-bit polynomial-time syntactic tag,
the relation is source-specific and still uses one retained bit and
polynomial preprocessing, without computing any compatibility matrix or SAT
witness.
\end{theorem}

\begin{proof}
For a fixed $a$, the intersection
\[
 \bigcap_{x\in X_a}R_h(x,\star)
 =\bigcap_{x\in X_a}\{h(x)\}
\]
is nonempty exactly when all displayed singleton sets are equal, which is
equivalent to constancy of $h$ on $X_a$.  Constancy on all fibers is exactly
factorization through $b$: define $f(a)$ to be the common value.  Conversely
$h=f\circ b$ is constant on each fiber.

If constancy fails, the compiler evaluates $h(x)$ before the cut, stores that
value, and emits it after the cut.  Its cost is precisely the declared cost
of $h$ plus the response storage/output cost.  No operation on $b$ is needed.
\end{proof}

\begin{corollary}[One-bit clean-core counterexample]
Let $F_0,F_1$ be two distinct equal-answer canonical clean-core formulas of
the same encoded length; the verified corpus supplies equal-size satisfiable
paired-incidence Monotone $3$SAT-$(2,2)$ examples.  Define
\[
 h(F)=\mathbf 1[F=F_0]
\]
on that length layer, defaulting to zero elsewhere.  Equality with the fixed
canonical string is decidable by one scan in $O(N)$ time and $O(\log N)$
workspace.  The singleton response sets for $F_0,F_1$ are disjoint although
$\SAT(F_0)=\SAT(F_1)$.  Thus empty same-answer response intersection can be
created by a trivial syntactic tag and is not a SAT lower bound.
\end{corollary}

The construction may be made even more representation-relative by choosing a
literal, order, padding, or wrapper tag.  Those variants strengthen the
hostile control; they are not invariant obstructions.

\section{Online absorption trilemma}

Combining the new theorem with FIBERDICH1 and previously verified
STATICNP1/PHABSORB1 gives the following exact synthesis boundary.

\begin{enumerate}[leftmargin=2em]
\item If every same-answer response-set intersection is nonempty and a common
choice is efficiently evaluable, the online obligation collapses to the SAT
answer bit.
\item If an intersection is empty but the separating response is a cheap
source function, source specificity is real and computational hardness is
absent.
\item If the response relation is total, polynomially balanced, and
fixed-level-PH verifiable, then under $P=NP$ prefix search constructs a valid
response in polynomial time.  This is the already verified STATICNP1 / PHABSORB1
closure, not a new theorem.
\end{enumerate}

Escaping all three branches requires a quantitative theorem, not merely a new
relation.  One must compile every assumed polynomial SAT decider into one
responder with fixed, solver-exponent-independent simultaneous resource
bounds and then place that same ordinary computation in a model with an
unconditional lower bound.

Therefore \claim{FIBER-SEPARATION-AS-OBSTRUCTION1} is \RETIRED.  The full
\claim{SOURCE-FIBER-SEPARATING-COMPILER1} from FIBERDICH1 is not contradicted
because it explicitly demanded a model-matched forbidden landing; its
fiber-separation subcondition is now classified as necessary interface
hygiene only.  The repaired team target is
\claim{QUANTITATIVE-TWO-WAY-ONLINE-TRANSFER1}, stated in Section~\ref{sec:target}.

\section{Finite reproducible audit}

The script
\path{experiments/review/source_specificity_no_hardness_synthesis.py}
exhausted all $16$ Boolean decision maps and $16$ Boolean source tags on four
sources.  Of the $256$ pairs, $60$ tags factored through the decision and
$196$ were source-specific.  Every one of the $16$ decision maps had a
same-answer collision and a one-bit equality tag separating such a pair.

It also exhausted all $256$ subsets of the eight proper three-variable
3-clauses.  There were $255$ satisfiable formulas and one unsatisfiable
formula.  Four syntactic tags separated same-answer sources; the three
balanced tags each separated $16{,}256$ same-answer pairs.  DIRECT, COPY,
DEAD, DELAY, INTERP, and RECOMP wrappers contributed $3{,}072$ checks;
independent-product and padding controls contributed $288$.  In total,
$4{,}677$ assertions passed in $0.024$ seconds with zero failures.

This evidence is \COMP.  The asymptotic conclusion comes only from the
symbolic theorem.  The raw output is
\path{results/raw/review/2026-10-12_source_specificity_no_hardness.json}.

\section{Newest specialist-cycle reconciliation}

\begin{itemize}
\item LB-0022 proves that exponentially many compatibility rows and
satisfying assignments on a canonical cubic Tseitin exact-$(4,4)$ family
collapse to the parity of the charges, computable in $O(N)$ time and
$O(\log N)$ workspace.  This is an exact easy-algebraic control, not the
active H22 hard-family theorem.
\item TESTS-0019 proves a polynomial-length unique bit-local pointer can be
validated and erased with logarithmic control space; validation time remains
charged.
\item PADQUOT1 proves full unpointed padding invariance is constant on the
satisfiable clean core; pointing avoids collapse by importing provenance.
\item FIBERDICH1 proves the exact common-response criterion.  The present run
shows that its opposite branch---source separation---also admits trivial
cheap tags and therefore is not a hardness certificate.
\end{itemize}

No specialist theorem supplies the missing arbitrary-decider transfer or
solver-exponent-independent endpoint.

\section{Actual main4/main5 comparison}

The current repository files were read directly.  \path{main4.tex} is still
the ARC achievable-region/source-pass entry point and \path{main5.tex} is
still the WRAPPER construction entry point.  Both include their historical
construction and recovery additions through MODELBRIDGE-0020 and explicitly
leave all-input acquisition, successor coverage, random access, replay, and
simultaneous resource closure open.  Neither file contains FIBERDICH1, the
present theorem, or an arbitrary-decider online compiler.  This report does
not silently reinterpret their partial domains as a universal theorem.

\section{SAME\_REALIZATION and lower-bound endpoint}

FLMV's directly transferable theorem is model-specific: for every constant
$c<\varphi$ there is $d>0$ such that SAT is not in simultaneous
$\mathrm{DTISP}(N^c,N^d)$ on the stated deterministic random-access machine
model.  Its locally computed reduction bits are useful only after one actual
machine has the required time and space.  It does not force a static SAT
decider to implement an online response relation.

Yao-style communication lower bounds transfer exactly only after the two
parties, input split, messages, and protocol are fixed.  Treating a
source-removal protocol as an arbitrary SAT machine is a \textbf{FAILED
TRANSFER}.  VLSI routing, streaming, observability, and proof-system strategy
extraction are \textbf{ANALOGY/CANDIDATE} unless an explicit simulation
places every SAT decider in the same model.

All compilation, source reads, validation, challenge generation, retained
state, rereading, recomputation, response output, conversion to ordinary SAT,
failed routes, time, and peak workspace must belong to one realization.

\section{Barrier and post-quantum audit}

The singleton theorem and the STATICNP1/PHABSORB1 closures relativize and
algebrize; they provide no escape from Baker--Gill--Solovay or
Aaronson--Wigderson.  No large constructive circuit property is used, so
Natural Proofs is not directly triggered or bypassed.

Collision entropy is an \textbf{EXACT OBJECT TRANSFER} for a declared
source-to-message map, but a low-collision syntactic tag can be easy; the
statistic does not lower-bound SAT decision time.  Special soundness is an
\textbf{EXACT THEOREM TRANSFER} only after one explicit interactive protocol
supplies two valid openings to a common commitment.  Its transfer to arbitrary
static SAT remains missing.  Schur squares, quadratic relation spaces,
Hilbert functions, syzygies, and cryptographic/random-oracle binding are
\textbf{FAILED TRANSFERS} here because no forced code or algebraic
representation occurs in the theorem.  LWE-style black-box extraction remains
an \textbf{ANALOGY/CANDIDATE} needing an unconditional worst-case transfer.

\section{Status ledger}

\begin{center}\small
\begin{tabular}{>{\raggedright\arraybackslash}p{0.47\linewidth}p{0.21\linewidth}p{0.24\linewidth}}
\toprule
Claim & Status & Scope\\ \midrule
\claim{SOURCE-SPECIFICITY-NO-HARDNESS1} & \EXACT & Singleton response relations; review pending\\
FIBERDICH1 common-response theorem & \EXACT & Deterministic total relation; review pending\\
STATICNP1 / PHABSORB1 response search & \COND & Under $P=NP$\\
Finite source-tag audit & \COMP & $4{,}677$ assertions\\
\claim{FIBER-SEPARATION-AS-OBSTRUCTION1} & \RETIRED & Source distinction alone\\
\claim{QUANTITATIVE-TWO-WAY-ONLINE-TRANSFER1} & \OPEN & All costs; same machine\\
Nonrelativizing/nonalgebrizing hinge & \OPEN & Arbitrary computation\\
$P\ne NP$ & \OPEN & No proof\\
\bottomrule
\end{tabular}
\end{center}

\section{Single highest-value target}\label{sec:target}

Prove or falsify \claim{QUANTITATIVE-TWO-WAY-ONLINE-TRANSFER1}.  It must give
one fixed full-language construction satisfying both directions:
\[
 \forall A\in\operatorname{Dec}_{\rm uniform}(\SAT)\cap\mathsf P,
 \quad A\mapsto P_A
\]
with explicit simultaneous resource exponents independent of the unknown
clock exponent of $A$, and
\[
 P_A\mapsto M_A\in\mathrm{DTISP}(N^c,N^d)
\]
for one FLMV-forbidden pair $(c,d)$, with all transformations and costs in the
same ordinary random-access realization.  The source-specific relation must
be decision-entangled in the precise sense supplied by the second map; empty
fiber intersection alone is insufficient.

This is separation-strength and remains \OPEN.

\section{CORPUS REVIEWED}

The run surveyed the complete GitHub tree at head
\texttt{b73fa9f2653e930c8406e896b1c8ba859fe25ca5}; read root
\path{main.tex}, \path{AGENTS.md}, \path{CODEX.md}, the pnp-research skill and
both protocol references; read every task under \path{tasks/}; read actual
\path{main4.tex} and \path{main5.tex}; read MAIN SYNTHESIS-0019; read newest
LB-0022, TESTS-0019, PADQUOT1, and FIBERDICH1 plus FORKFACTOR1; and reconciled
the relevant HEC, VSC-HEC, G-PPC-HEC, CRS-PPC, CCSPC, STATICNP1, PHABSORB1,
ROUNDABSORB1, GLOBALGAME1, SOURCESTOP1, and VIEWSEP1 lineage.  Governance,
cited experiments/results, and repository searches for fiber, source-specific
relation, response absorption, and STATICNP1 were inspected.

Primary literature checked before the theorem includes FLMV's original
time--space paper, Yao's fixed-protocol communication framework,
Baker--Gill--Solovay, Aaronson--Wigderson, and Razborov--Rudich.

\section{DUPLICATION / CONTRADICTION AUDIT}

STATICNP1 already proves NP-level continuation search is polynomial under
$P=NP$; PHABSORB1 and ROUNDABSORB1 extend that closure.  The present report
does not rename those results.  VIEWSEP1 already establishes separating-view
necessity for incompatible responses.  The new exact contribution is its
converse warning: separation can be a trivial one-bit syntactic tag and does
not itself imply hardness.  No contradiction with FIBERDICH1 occurs; the two
theorems classify opposite sides of the same intersection test.

\section{WHAT CHANGED RELATIVE TO ROOT MAIN}

Root MAIN previously ended at the pointed wrapper/padding target and had not
yet synthesized LB-0022, TESTS-0019, or FIBERDICH1.  The current reconciliation
adds the singleton-relation counterexample, retires fiber separation as a
standalone obstruction, and replaces the executive target by a quantitative
two-way SAME\_REALIZATION transfer.  The exact H22 clean core and unrestricted
runtime target are unchanged.

\section{ADOPTED / REJECTED / RETIRED PRIOR COMPONENTS}

\textbf{Adopted:} FIBERDICH1, STATICNP1, PHABSORB1, VIEWSEP1, PADQUOT1,
FORKFACTOR1, TESTS-0019 pointer erasure, LB-0022 as an easy affine control,
SAME\_REALIZATION, and FLMV only in its declared model.

\textbf{Rejected:} empty same-answer intersection, source-specific tags,
pointer identity, collision count, or protocol communication as an automatic
arbitrary-SAT runtime lower bound.

\textbf{FALSIFIED/RETIRED:} \claim{FIBER-SEPARATION-AS-OBSTRUCTION1},
answer-derived superconstant state, unpointed padding-normalized invariants,
and source tags as hardness resources without a quantitative endpoint map.

\section{STALE OR OVERSTATED CLAIMS FOUND}

Any claim that SOURCE-FIBER-SEPARATING-COMPILER1 has materially advanced
merely because the response sets distinguish two same-answer sources is
overstated.  The compiler target remains meaningful only with its previously
stated all-cost, model-matched forbidden endpoint.  Likewise, the older
main4/main5 partial construction points are not universal SAT obligations.

\begin{thebibliography}{9}
\bibitem{FLMV05} L. Fortnow, R. Lipton, D. van Melkebeek, and A. Viglas,
``Time-space lower bounds for satisfiability,'' \emph{J. ACM} 52(6), 2005.
\url{https://doi.org/10.1145/1101821.1101828}
\bibitem{Yao79} A. C.-C. Yao, ``Some complexity questions related to
distributive computing,'' \emph{STOC 1979}, 209--213.
\url{https://doi.org/10.1145/800135.804414}
\bibitem{BGS75} T. Baker, J. Gill, and R. Solovay,
``Relativizations of the P=?NP question,'' \emph{SIAM J. Comput.} 4(4), 1975.
\url{https://doi.org/10.1137/0204037}
\bibitem{AW09} S. Aaronson and A. Wigderson,
``Algebrization: a new barrier in complexity theory,'' \emph{ACM TOCT} 1(1), 2009.
\url{https://eccc.weizmann.ac.il/report/2008/034/}
\bibitem{RR97} A. Razborov and S. Rudich, ``Natural proofs,''
\emph{J. Comput. Syst. Sci.} 55(1), 1997.
\url{https://doi.org/10.1006/jcss.1997.1494}
\end{thebibliography}

% END_OF_VERIFIED_MAIN::SYNTHESIS-0020
\end{document}
